%% 2. 手法（骨組みの 3・5・6）
\setband{\T{\S2 手法と比較}{\S2 Methods and comparisons}}

\begin{frame}{\T{既存研究の不足：「上がるか」だけを測り，どのモダリティで・損のどの部分が減るかを測らない}{What existing work misses: it measures ``does it help'' but not for which modality or which part of the loss}}
\rl{\T{既存研究（骨組み 3）}{Prior work (skeleton 3)}}
\centering
\begin{tikzpicture}[font=\scriptsize, x=1cm, y=1cm]
  % (a)
  \node[qbox, text width=3.9cm, minimum height=2.5cm, anchor=north] (a) at (0,0) {\textbf{(a) \T{同じモダリティのラベル無しの手法}{Unlabeled-data methods within one modality}}\\[2pt]
    \T{自己学習，エントロピー最小化，T3A，ラベル伝播，DANN，CORAL，MME …}{self-training, entropy min., T3A, label propagation, DANN, CORAL, MME, …}\\[3pt]
    \T{疑似ラベルは\textbf{生徒自身}の予測か，特徴の分布合わせ．生徒より当たる予測は入らない}{Pseudo-labels come from the \textbf{student itself}, or features are aligned. No prediction better than the student enters}};
  % (b)
  \node[qbox, text width=3.9cm, minimum height=2.5cm, anchor=north] (b) at (4.6,0) {\textbf{(b) \T{別のモダリティ 1 つの予測}{One other modality's prediction}}\\[2pt]
    \T{Albanie ほか 2018：顔の識別器の予測で音声の識別器を学ぶ}{Albanie et al.\ 2018: the face classifier teaches the audio classifier}\\[3pt]
    \T{1 つのモダリティの予測は弱いことがあり，学ぶ側を悪くしうる（TelME 2024：言葉が音声から学ぶと 66.6 → 57.4）}{A single modality can be weak and hurt the learner (TelME 2024: text learning from audio, 66.6 → 57.4)}};
  % (c)
  \node[propbox, text width=3.9cm, minimum height=2.5cm, anchor=north] (c) at (9.2,0) {\textbf{(c) \T{この研究}{This work}}\\[2pt]
    \T{3 モダリティの識別器を学び直し，確率の\textbf{積}を疑似ラベルにする．K 本と出発点への引き戻しを足す}{Retrain the 3 unimodal classifiers and use the \textbf{product} of their probabilities as pseudo-labels, plus K labels and a pull toward the start}\\[3pt]
    \T{そして\textbf{測る}：どのモダリティの生徒で，損のどの部分（切片か重みの向きか）が減るか}{And \textbf{measure}: for which student, and which part of the loss (intercept or weight direction) shrinks}};
\end{tikzpicture}

\vspace{0.4em}
\takeaway{\T{骨組み 3 の文：既存の評価は，ほかのモダリティの予測で学ぶと成績が上がるかだけを測り，どのモダリティで，損のどの部分（切片か重みの向きか）が減るかを測らない}{Skeleton item 3: existing evaluations only ask whether learning from other modalities' predictions helps, not for which modality or which part of the loss (intercept or weight direction) shrinks}}
\end{frame}

\begin{frame}{\T{提案の流れ：ソースの識別器を出発点にし，3 モダリティの積の予測と K 本で生徒を学ぶ}{Proposal: start from the source classifiers; train the student on K labels plus the product of three modalities' predictions}}
\rl{\T{手法}{Method}}
\hosoku{appsolver}{\T{学習器と引き戻しの詳細}{solver and pull details}}
\centering
\begin{tikzpicture}[font=\scriptsize, x=1cm, y=1cm]
  % 行 A：出発点
  \node[srcbox, text width=3.6cm, minimum height=0.9cm] (src) at (0,0) {\T{ソースのコーパスごとの識別器}{Per-source classifiers}\\ \T{（3 コーパス × A・V・T・AVT）}{(3 corpora × A, V, T, AVT)}};
  \node[clf, text width=2.6cm, minimum height=0.9cm] (start) at (4.3,0) {\T{出発点}{Start}\\ \T{＝重みと切片の平均}{= mean of weights and intercepts}};
  \draw[arr] (src) -- node[above, font=\tiny]{\T{平均}{average}} (start);
  \node[font=\tiny, color=mMuted, anchor=west] at (6.0,0) {\T{すべての手法が同じ出発点を使う}{every method uses the same start}};
  % 行 B：教師
  \node[unlbox, text width=2.9cm, minimum height=1.0cm] (avtk) at (0.35,-1.7) {\T{出発点の AVT を}{AVT start, retrained}\\ \T{K 本で学び直す → 候補に予測}{with K labels → predict the pool}};
  \node[unlbox, text width=3.3cm, minimum height=1.0cm] (ret) at (4.3,-1.7) {\T{A・V・T の識別器を出発点から}{Retrain A, V, T from the start on}\\ \T{K 本 ＋ その予測で学び直す}{K labels + those predictions}};
  \node[modbox, text width=3.3cm, minimum height=1.0cm] (q) at (8.6,-1.7) {\T{\textbf{教師}：3 つの確率の積}{\textbf{Teacher}: product of the 3 probabilities}\\ \T{→ 候補の各発話に疑似ラベル Q}{→ pseudo-label Q for each pool utterance}};
  \draw[arr] (avtk) -- (ret); \draw[arr] (ret) -- node[above, font=\tiny]{\T{掛ける}{multiply}} (q);
  \draw[arr] (start.south) -- ++(0,-0.35) -| (ret.north);
  % 行 C：生徒
  \node[propbox, text width=8.6cm, minimum height=1.25cm] (stu) at (4.3,-3.6) {\T{\textbf{生徒}（A／V／T／AVT）を出発点から学ぶ．損失 ＝}{\textbf{Student} (A / V / T / AVT), trained from the start. Loss =}\\
    \colorbox{cTgt}{\T{K 本の正解との交差エントロピー}{cross-entropy on the K labels}} ＋
    \colorbox{cMod}{\T{疑似ラベル Q との交差エントロピー（感情ごとに総量をそろえる）}{cross-entropy to Q (total weight balanced per emotion)}} ＋
    \colorbox{cSrc}{\T{出発点への引き戻し}{pull toward the start}}};
  \draw[arr] (q.south) |- ++(-1.5,-0.55) -| (stu.north east);
  \node[lblbox, text width=1.6cm] (k) at (-1.1,-3.6) {K \T{本}{labels}};
  \draw[arr] (k) -- (stu);
  \node[testbox, text width=1.4cm] (test) at (10.5,-3.6) {\T{テストで UAR}{UAR on test}};
  \draw[arr] (stu) -- (test);
\end{tikzpicture}
\notreal
\end{frame}

\begin{frame}{\T{提案と温度つき蒸留の違いは 2 点だけで，評価ではほぼ同じ値になった}{The proposal differs from temperature distillation in only two points, and scored the same in evaluation}}
\rl{\T{手法}{Method}}
\begin{columns}[T]
\begin{column}{0.58\textwidth}
\small
\begin{tabular}{@{}p{0.30\linewidth}p{0.31\linewidth}p{0.31\linewidth}@{}}
\toprule
& \T{提案}{Proposal} & \T{温度つき蒸留（Hinton 2015，Albanie 2018 の形）}{Temperature distillation (Hinton 2015; Albanie 2018 form)} \\
\midrule
\T{教師}{Teacher} & \T{同じ Q}{same Q} & \T{同じ Q}{same Q} \\
\T{疑似ラベルの形}{Pseudo-label form} & \T{Q そのまま}{Q as is} & \T{Q を温度 2 で柔らげる}{Q softened with temperature 2} \\
\T{感情ごとの総量}{Per-emotion total weight} & \T{そろえる}{balanced} & \T{そろえない}{not balanced} \\
\T{K 本・引き戻し}{K labels, pull} & \T{同じ}{same} & \T{同じ}{same} \\
\T{つまみ（格子）}{Knobs (grid)} & \T{全体の重み ρ，K 本の倍率}{overall weight ρ, label multiplier} & \T{温度，全体の重み ρ}{temperature, overall weight ρ} \\
\bottomrule
\end{tabular}
\end{column}
\begin{column}{0.42\textwidth}
\small
\T{%
\begin{itemize}
\item 評価の回で，提案 − 蒸留 は 48 セル中 40 セルで区間が 0 を含む（\S3）
\item 残りの差も −1.2〜+2.2 ポイントで，向きはコーパスで分かれる
\item したがって解析では\textbf{両者を区別せず}，「別のモダリティの予測（教師）で学ぶ手法」として扱う
\end{itemize}
}{%
\begin{itemize}
\item In evaluation, proposal − distillation has an interval containing 0 in 40 of 48 cells (\S3)
\item The remaining differences are −1.2 to +2.2 points, with signs varying by corpus
\item Hence the analysis \textbf{does not distinguish} them: both are ``learning from the other modalities' predictions (the teacher)''
\end{itemize}
}
\end{column}
\end{columns}
\takeaway{\T{「蒸留が良い理由」を問うときの「蒸留」は，この 2 手法の共通部分＝教師 Q で学ぶことを指す}{When we ask ``why distillation helps'', ``distillation'' means what the two share: learning from the teacher Q}}
\end{frame}

\begin{frame}{\T{評価の尺度と比較：UAR で測り，使う情報ごとに相手を分ける}{Measure and comparisons: UAR, with baselines grouped by the information they use}}
\rl{\T{評価（骨組み 5・6）}{Evaluation (skeleton 5, 6)}}
\hosoku{appmethods}{\T{手法と格子の一覧}{list of methods and grids}}
\begin{columns}[T]
\begin{column}{0.46\textwidth}
\footnotesize
\textbf{\T{尺度（骨組み 5）}{Measure (skeleton 5)}}
\T{%
\begin{itemize}
\item UAR（4 感情の再現率の平均）．MELD と MOSEI は Neutral が半分近くを占め，正解率だと「全部 Neutral」が高く出るため
\item テストの話者は学習側に入れない（MOSEI は話者の代わりにチャンネルか動画）
\end{itemize}
\textbf{除去（骨組み 6）}：提案から 1 つずつ外した版を実験 1 で計算済み
\begin{itemize}
\item (i) 教師を学び直さない　(ii) 教師を 1 モダリティにする
\item (iii) 総量をそろえない（＝温度 1 の蒸留）　(iv) 切片だけを K 本で学ぶ
\end{itemize}
}{%
\begin{itemize}
\item UAR (mean recall over 4 emotions). In MELD and MOSEI nearly half the utterances are Neutral, so accuracy rewards ``always Neutral''
\item Test speakers never enter training (MOSEI: channel or video stands in for speaker)
\end{itemize}
\textbf{Ablations (skeleton 6)}: versions with one part removed, already computed in Exp.~1
\begin{itemize}
\item (i) teacher not retrained\quad (ii) single-modality teacher
\item (iii) no balancing (= distillation at T = 1)\quad (iv) intercept only from K labels
\end{itemize}
}
\end{column}
\begin{column}{0.54\textwidth}
\footnotesize
\textbf{\T{比較の相手（骨組み 6）}{Baselines (skeleton 6)}}
\begin{tabular}{@{}p{0.30\linewidth}p{0.66\linewidth}@{}}
\toprule
\T{使う情報}{Information used} & \T{相手}{Baselines} \\
\midrule
\T{K 本だけ}{K labels only} & \T{比較の基準（出発点から K 本で学び直す）}{the baseline (start retrained on K labels)} \\
\T{K 本 ＋ 同じ教師 Q}{K labels + the same teacher Q} & \T{温度つき蒸留，一致した発話だけ，自分の予測と混ぜる，柔らかい共訓練，共正則化（実験 1）}{temperature distillation, agreed-only, mixing with own prediction, soft co-training, co-regularization (Exp.~1)} \\
\T{K 本 ＋ ラベル無し（同じモダリティ）}{K labels + unlabeled (same modality)} & \T{自己学習，エントロピー最小化，IM，T3A，ラベル伝播，自分の予測を提案の形で（実験 2）}{self-training, entropy min., IM, T3A, label propagation, own prediction in the proposal's form (Exp.~2)} \\
\T{＋ ソースのデータ}{+ source data} & \T{S+T，CORAL，DANN，MME，xMUDA 型（実験 2）}{S+T, CORAL, DANN, MME, xMUDA-style (Exp.~2)} \\
\bottomrule
\end{tabular}
\vspace{0.3em}
{\scriptsize\color{mMuted}\T{全手法が同じ出発点・同じ標準化・同じ K 本を使い，つまみは仮のターゲットで選んで固定した}{All methods share the start, the standardization and the K labels; knobs were selected on proxy targets and frozen}}
\end{column}
\end{columns}
\end{frame}

\begin{frame}{\T{統計の決まり：検定に残るのは 128 セル中 15 セルで，残りは区間つきの記述}{Statistical rules: 15 of 128 cells are tested; the rest are descriptions with intervals}}
\rl{\T{統計}{Statistics}}
\hosoku{appstats}{\T{区間と検定の規則の詳細}{details of intervals and tests}}
\small
\T{%
\begin{itemize}
\item セル ＝（コーパス，生徒，K）．族 1（提案 − K 本だけ），族 3（提案 − 同じ情報の相手）など，コーパスごとに Holm で補正する
\item 区間は，テストの話者の塊の引き直し（jackknife）と K 本の抜き方の揺れを含む．AFEW・IEMOCAP・MELD は実効の塊が 15・5・9 と少なく，1〜2 ポイントの差は有意にならない
\item そこで評価の\textbf{前}に，検出力の見込みが足りるセル（15）だけを検定に残し，残り 113 セルは「区間つきの記述」と決めた．記述のセルでは「差が無い」「効かない」と書かない
\item 手法の形は評価と同じ 4 コーパスの上で作ったので，p 値は確認の検定ではない（論文の限界に書く）
\end{itemize}
}{%
\begin{itemize}
\item A cell = (corpus, student, K). Family 1 (proposal − K-only), family 3 (proposal − same-information baseline), etc., with Holm correction per corpus
\item Intervals include resampling of test speaker clusters (jackknife) and the random draw of the K labels. AFEW, IEMOCAP and MELD have only 15, 5 and 9 effective clusters, so 1--2 point differences cannot be significant
\item Therefore, \textbf{before} evaluation, only the 15 cells with sufficient expected power were kept for testing; the other 113 are ``descriptions with intervals''. For those we never write ``no difference'' or ``does not help''
\item Method forms were built on the same 4 corpora used for evaluation, so p-values are not confirmatory (stated in the limitations)
\end{itemize}
}
\end{frame}
